% Môn: Vật lý
% Lớp: 11
% Chương: Chương I. Dao động
% Bài: L11C1 Bài 5. Động năng. Thế năng. Sự chuyển hóa giữa động năng và thế năng trong dao động điều hòa · Đơn vị và đại lượng trong dao động điều hòa
% ===== Câu 29 =====
% ID: L11C1B5-105-TN
% Mức: VD
\dangbt{Đơn vị và đại lượng trong dao động điều hòa}
\begin{ex}
% ID: L11C1B5-105-TN
% Mức: VD
Số dao động toàn phần vật thực hiện trong một giây là
\choice
{\True tần số}
{chu kì}
{tần số góc}
{pha dao động}
\end{ex}

% ===== Câu 107 =====
% ID: L11C1B5-106-DS
% Mức: TH
\begin{ex}
% ID: L11C1B5-106-DS
% Mức: TH
Hãy xác định tính đúng sai của các phát biểu sau về dao động cơ và dao động điều hòa.
\choiceTF
{\True Dao động cơ là chuyển động qua lại quanh một vị trí cân bằng.}
{\True Dao động điều hòa có li độ biến thiên theo hàm sin hoặc cos của thời gian.}
{Chu kì của mọi dao động tự do luôn phụ thuộc vào biên độ dao động.}
{\True Pha dao động cho phép xác định trạng thái dao động tại một thời điểm.}
\loigiai{
\begin{itemchoice}
\item (Đúng) Dao động cơ là chuyển động qua lại quanh một vị trí cân bằng.
\item (Đúng) Dao động điều hòa có li độ biến thiên theo hàm sin hoặc cos của thời gian.
\item (Sai) Chu kì của mọi dao động tự do luôn phụ thuộc vào biên độ dao động.
\item (Đúng) Pha dao động cho phép xác định trạng thái dao động tại một thời điểm. (Vì): ao động cơ diễn ra quanh vị trí cân bằng; dao động điều hòa có dạng sin/cos; chu kì dao động tự do phụ thuộc đặc tính hệ và pha xác định trạng thái dao động
\end{itemchoice}
}
\end{ex}

% ===== Câu 109 =====
% ID: L11C1B5-107-DS
% Mức: TH
\begin{ex}
% ID: L11C1B5-107-DS
% Mức: TH
Hãy xác định tính đúng sai của các phát biểu sau về các đại lượng đặc trưng của dao động điều hòa.
\choiceTF
{\True Biên độ là đại lượng không âm.}
{\True Đơn vị của tần số góc là rad/s.}
{\True Tần số bằng nghịch đảo của chu kì.}
{Pha ban đầu không phụ thuộc vào cách chọn gốc thời gian.}
\loigiai{
\begin{itemchoice}
\item (Đúng) Biên độ là đại lượng không âm.
\item (Đúng) Đơn vị của tần số góc là rad/s.
\item (Đúng) Tần số bằng nghịch đảo của chu kì.
\item (Sai) Pha ban đầu không phụ thuộc vào cách chọn gốc thời gian.
\end{itemchoice}
}
\end{ex}

% ===== Câu 133 =====
% ID: L11C1B5-108-TN
% Mức: TH
\begin{ex}
% ID: L11C1B5-108-TN
% Mức: TH
Bài 1 – đơn vị và đại lượng: Đại lượng nào có đơn vị rad/s?
\choice
{\True Tần số góc}
{Chu kì}
{Biên độ}
{Tần số}
\end{ex}

% ===== Câu 134 =====
% ID: L11C1B5-109-TN
% Mức: VD
\begin{ex}
% ID: L11C1B5-109-TN
% Mức: VD
Bài 1 – đơn vị và đại lượng: Đại lượng nào biểu thị độ lệch lớn nhất khỏi vị trí cân bằng?
\choice
{\True Biên độ}
{Pha ban đầu}
{Chu kì}
{Tần số}
\end{ex}
